\chapter{Linked Lists}
\chapterepigraph{A linked list trades the array's random access for
something arrays cannot do cheaply: rearranging elements by rewiring
pointers instead of shifting memory. Almost every linked-list problem is
won or lost on two techniques -- the two-pointer (slow/fast) trick and the
disciplined manipulation of \texttt{next} pointers -- plus the habit of
using a dummy head node so the ``first element'' is never a special case.}

\section{Pointers, Two Pointers, and the Dummy Node}
\label{sec:ll-intro}

A node holds a value and a pointer to the next node; the list is a chain
ending in \texttt{nullptr}. There is no index, so everything is done by
walking pointers -- which is why a handful of pointer patterns solve the
whole chapter.

\begin{ideabox}[Three Habits That Solve Most List Problems]
\begin{itemize}
  \item \textbf{Slow/fast (tortoise--hare):} advance one pointer by one and
        another by two. This finds the middle, detects a cycle, and locates
        the $n$-th node from the end in a single pass.
  \item \textbf{Rewire, don't copy:} to reverse, reorder, or remove nodes,
        change \texttt{next} pointers; keep a \texttt{prev} and a
        \texttt{next} handle so you never lose the rest of the list.
  \item \textbf{Dummy head:} allocate a throwaway node before the real head
        so insertion/deletion at the front is not a special case -- return
        \texttt{dummy.next} at the end.
\end{itemize}
\end{ideabox}

\subsection*{How to Recognise the Technique}

\begin{patternbox}[Linked-List Keywords]
\begin{itemize}
  \item \textbf{``Middle,'' ``cycle,'' ``$n$-th from end,'' ``loop
        length''} -- slow/fast two pointers.
  \item \textbf{``Reverse,'' ``reorder,'' ``swap in groups,'' ``rotate''}
        -- pointer rewiring, often with a dummy node.
  \item \textbf{``Palindrome,'' ``add two numbers,'' ``intersection''} --
        reverse or align sub-lists, then walk in step.
  \item \textbf{``Sort,'' ``0s/1s/2s''} -- merge sort on the list, or
        pointer partitioning.
  \item \textbf{``Flatten,'' ``clone with random pointer''} -- recursion or
        an interleaving trick that reuses the list's own structure.
  \item \textbf{Doubly linked list} -- an extra \texttt{prev} pointer makes
        backward traversal and $O(1)$ deletion possible.
\end{itemize}
\end{patternbox}

Throughout, a singly linked node is
\texttt{struct ListNode \{ int val; ListNode* next; \}} and a doubly linked
node additionally carries \texttt{ListNode* prev}. Every listing is a
complete, compilable program that builds a small list, runs the operation,
and prints the result.

\newpage

\section{Singly Linked List: Construction and Basic Operations}
\label{sec:ll-singly-basics}

\problemstatement{Implement the foundational operations on a singly linked
list: build from an array, traverse, find the length, search for a value,
insert at the head or tail, and delete a value.}

\begin{patternbox}
No pattern recognition is needed here -- this is the machinery every later
section rests on. The one habit to internalise is keeping a \texttt{prev}
handle while walking, so a node can be unlinked or a new node spliced in
without losing the rest of the chain.
\end{patternbox}

\subsection*{Walking with a pointer}

Each node stores a value and a pointer to the next node; the last node
points to \texttt{nullptr}. Traversal is a loop that follows \texttt{next}
until null. Insertion at the head is $O(1)$ (point the new node at the old
head). Insertion at the tail and deletion require walking to the right spot,
keeping the previous node so its \texttt{next} can be rewired.

\begin{ideabox}[Keep the Previous Node When You Might Rewire]
Deletion and mid-list insertion both need to change the \texttt{next} of the
node \emph{before} the target. Tracking \texttt{prev} as you traverse makes
these $O(1)$ once the position is found, with no re-walk.
\end{ideabox}

\subsection*{C++ implementation}

\begin{lstlisting}
#include <bits/stdc++.h>
using namespace std;

struct ListNode {
    int val;
    ListNode* next;
    ListNode(int v) : val(v), next(nullptr) {}
};

ListNode* build(const vector<int>& a) {           // build from array
    ListNode dummy(0), *tail = &dummy;
    for (int x : a) { tail->next = new ListNode(x); tail = tail->next; }
    return dummy.next;
}

int length(ListNode* head) {
    int n = 0;
    for (ListNode* c = head; c; c = c->next) n++;
    return n;
}

bool search(ListNode* head, int key) {
    for (ListNode* c = head; c; c = c->next) if (c->val == key) return true;
    return false;
}

ListNode* insertHead(ListNode* head, int v) {
    ListNode* node = new ListNode(v);
    node->next = head;
    return node;                                  // new head
}

ListNode* deleteValue(ListNode* head, int key) {
    ListNode dummy(0); dummy.next = head;
    ListNode* prev = &dummy;
    while (prev->next && prev->next->val != key) prev = prev->next;
    if (prev->next) prev->next = prev->next->next;   // unlink
    return dummy.next;
}

void print(ListNode* head) {
    for (ListNode* c = head; c; c = c->next) cout << c->val << " ";
    cout << "\n";
}

int main() {
    ListNode* head = build({1, 2, 3, 4});
    print(head);                                  // 1 2 3 4
    cout << "length " << length(head) << ", search 3? " << search(head, 3) << "\n";
    head = insertHead(head, 0);
    head = deleteValue(head, 3);
    print(head);                                  // 0 1 2 4
    return 0;
}
\end{lstlisting}

\begin{complexitybox}
\textbf{Time Complexity.} $O(1)$ head insert; $O(n)$ for length, search,
tail insert, and delete. \textbf{Space Complexity.} $O(1)$ extra (plus the
list itself).
\end{complexitybox}

\begin{takeawaybox}
The dummy node plus a \texttt{prev} handle removes the ``deleting the head''
special case and is the scaffolding for nearly every operation in this
chapter. Master this template before the pointer-heavy problems ahead.
\end{takeawaybox}

\section{Doubly Linked List: Construction and Operations}
\label{sec:ll-doubly-basics}

\problemstatement{Implement a doubly linked list (DLL): build it, insert
before or after a node, and delete a node -- all in $O(1)$ once the node is
in hand, thanks to the backward \texttt{prev} pointer.}

\begin{patternbox}
A DLL adds a \texttt{prev} pointer to each node. The pay-off is $O(1)$
deletion of a known node (no need to walk to find its predecessor) and easy
backward traversal -- which is exactly why LRU caches and browser-history
structures use it.
\end{patternbox}

\subsection*{Two pointers per node}

Each node stores \texttt{prev} and \texttt{next}. Inserting a node $x$
between $a$ and $b$ rewires four pointers: $a.\texttt{next}=x$,
$x.\texttt{prev}=a$, $x.\texttt{next}=b$, $b.\texttt{prev}=x$. Deleting a
node $x$ splices its neighbours together: $x.\texttt{prev}.\texttt{next}=
x.\texttt{next}$ and $x.\texttt{next}.\texttt{prev}=x.\texttt{prev}$, with
null checks at the ends. Because a node knows its own predecessor, deletion
needs only the node itself.

\begin{ideabox}[The prev Pointer Buys O(1) Deletion]
In a singly linked list, deleting a node requires its predecessor, found by
an $O(n)$ walk. The \texttt{prev} pointer stores that predecessor, so any
node can be removed in constant time -- the feature that makes DLLs the
backbone of caches.
\end{ideabox}

\subsection*{C++ implementation}

\begin{lstlisting}
#include <bits/stdc++.h>
using namespace std;

struct Node {
    int val;
    Node* prev;
    Node* next;
    Node(int v) : val(v), prev(nullptr), next(nullptr) {}
};

Node* build(const vector<int>& a) {
    Node dummy(0), *tail = &dummy;
    for (int x : a) {
        Node* node = new Node(x);
        tail->next = node; node->prev = tail;
        tail = node;
    }
    Node* head = dummy.next;
    if (head) head->prev = nullptr;
    return head;
}

Node* deleteNode(Node* head, Node* x) {           // O(1) unlink
    if (x->prev) x->prev->next = x->next;
    else head = x->next;                          // x was the head
    if (x->next) x->next->prev = x->prev;
    return head;
}

void insertAfter(Node* a, int v) {                // splice x between a and a->next
    Node* x = new Node(v);
    x->next = a->next; x->prev = a;
    if (a->next) a->next->prev = x;
    a->next = x;
}

void print(Node* head) {
    for (Node* c = head; c; c = c->next) cout << c->val << " ";
    cout << "\n";
}

int main() {
    Node* head = build({1, 2, 3});
    insertAfter(head, 9);                          // 1 9 2 3
    print(head);
    head = deleteNode(head, head->next);           // remove the 9
    print(head);                                   // 1 2 3
    return 0;
}
\end{lstlisting}

\begin{complexitybox}
\textbf{Time Complexity.} $O(1)$ per insert/delete of a known node; $O(n)$
to build. \textbf{Space Complexity.} $O(1)$ extra per operation.
\end{complexitybox}

\begin{takeawaybox}
A doubly linked list is a singly linked list with a backward pointer,
trading a little memory for $O(1)$ deletion and two-way traversal. Keep the
four-pointer insert and two-pointer delete rewiring precise -- off-by-one
null handling at the ends is the usual bug.
\end{takeawaybox}

\section{Reverse a Doubly Linked List}
\label{sec:ll-reverse-dll}

\problemstatement{Reverse a doubly linked list in place so its head becomes
the tail and vice versa, returning the new head.}

\begin{patternbox}
Because each DLL node already knows both neighbours, reversing is just
\textbf{swapping every node's \texttt{prev} and \texttt{next}} pointers.
After swapping all of them, the old tail -- whose \texttt{next} is now null
-- is the new head.
\end{patternbox}

\subsection*{Swap the two pointers at every node}

Walk the list; at each node exchange its \texttt{prev} and \texttt{next}.
Once swapped, to advance you must follow the \emph{old} \texttt{next}, which
is now stored in \texttt{prev} -- so move via the just-swapped
\texttt{prev}. When you fall off the end, the last node you swapped is the
new head. No new nodes are created; only pointers are flipped.

\begin{ideabox}[Reversal Is Local Pointer Swapping]
In a DLL, the global reversal decomposes into an identical local operation
at each node: swap its two links. Because the links are symmetric, doing
this everywhere reverses the whole list without tracking a running
\texttt{prev} as a singly linked reversal must.
\end{ideabox}

\subsection*{C++ implementation}

\begin{lstlisting}
#include <bits/stdc++.h>
using namespace std;

struct Node {
    int val;
    Node* prev;
    Node* next;
    Node(int v) : val(v), prev(nullptr), next(nullptr) {}
};

Node* build(const vector<int>& a) {
    Node dummy(0), *tail = &dummy;
    for (int x : a) { Node* n = new Node(x); tail->next = n; n->prev = tail; tail = n; }
    Node* head = dummy.next; if (head) head->prev = nullptr;
    return head;
}

Node* reverse(Node* head) {
    Node* curr = head;
    Node* last = head;
    while (curr) {
        last = curr;
        swap(curr->prev, curr->next);             // flip both links
        curr = curr->prev;                        // old next is now in prev
    }
    return last;                                  // old tail = new head
}

void print(Node* head) {
    for (Node* c = head; c; c = c->next) cout << c->val << " ";
    cout << "\n";
}

int main() {
    Node* head = build({1, 2, 3, 4});
    head = reverse(head);
    print(head);                                  // 4 3 2 1
    return 0;
}
\end{lstlisting}

\begin{complexitybox}
\textbf{Time Complexity.} $O(n)$ -- one pass. \textbf{Space Complexity.}
$O(1)$ -- only pointer swaps.
\end{complexitybox}

\begin{takeawaybox}
Reversing a DLL is swapping \texttt{prev} and \texttt{next} at every node --
the symmetric structure makes global reversal a local operation. Contrast
with the singly linked reversal (next section) which must thread a running
\texttt{prev} because there is no backward pointer.
\end{takeawaybox}

\section{Middle of a Linked List}
\label{sec:ll-middle}

\problemstatement{Given the head of a singly linked list, return its middle
node in a single pass. For even length, return the second of the two middle
nodes.}

\begin{patternbox}
\emph{``Find the middle in one pass''} is the \textbf{slow/fast two-pointer}
(tortoise--hare) technique: move one pointer one step and another two steps;
when the fast pointer reaches the end, the slow pointer is at the middle.
\end{patternbox}

\subsection*{Two speeds, one pass}

Start both pointers at the head. Each iteration advances \texttt{slow} by
one node and \texttt{fast} by two. When \texttt{fast} (or
\texttt{fast->next}) runs off the end, \texttt{slow} has travelled exactly
half as far -- it sits at the middle. This avoids first computing the length
and then re-walking, doing everything in one traversal.

\begin{ideabox}[Half Speed Lands on the Middle]
Because \texttt{fast} covers twice the distance of \texttt{slow}, by the
time \texttt{fast} finishes the list \texttt{slow} has covered half of it.
The loop condition (\texttt{fast \&\& fast->next}) chooses the second middle
for even lengths.
\end{ideabox}

\subsection*{Algorithm}

\begin{enumerate}
  \item \texttt{slow = fast = head}.
  \item While \texttt{fast} and \texttt{fast->next} are non-null: advance
        \texttt{slow} by one, \texttt{fast} by two.
  \item Return \texttt{slow}.
\end{enumerate}

\subsection*{Dry run}

\dryrun{List $1\to2\to3\to4\to5$.}

\begin{itemize}
  \item slow$=1$,fast$=1$. Step: slow$=2$,fast$=3$. Step: slow$=3$,
        fast$=5$.
  \item fast->next is null, stop. Middle $=3$.
\end{itemize}

\subsection*{C++ implementation}

\begin{lstlisting}
#include <bits/stdc++.h>
using namespace std;

struct ListNode {
    int val; ListNode* next;
    ListNode(int v) : val(v), next(nullptr) {}
};

ListNode* build(const vector<int>& a) {
    ListNode dummy(0), *t = &dummy;
    for (int x : a) { t->next = new ListNode(x); t = t->next; }
    return dummy.next;
}

ListNode* middleNode(ListNode* head) {
    ListNode* slow = head;
    ListNode* fast = head;
    while (fast && fast->next) {
        slow = slow->next;
        fast = fast->next->next;
    }
    return slow;
}

int main() {
    ListNode* head = build({1, 2, 3, 4, 5});
    cout << middleNode(head)->val << "\n";        // 3
    ListNode* even = build({1, 2, 3, 4});
    cout << middleNode(even)->val << "\n";         // 3 (second middle)
    return 0;
}
\end{lstlisting}

\begin{complexitybox}
\textbf{Time Complexity.} $O(n)$ -- one pass. \textbf{Space Complexity.}
$O(1)$.
\end{complexitybox}

\begin{takeawaybox}
The slow/fast two-pointer finds the middle in one pass and is the single
most reused linked-list trick: cycle detection, $n$-th-from-end, and
palindrome checks all start from it. Learn its even/odd termination
precisely.
\end{takeawaybox}

\section{Reverse a Linked List (Iterative)}
\label{sec:ll-reverse-iterative}

\problemstatement{Reverse a singly linked list in place and return the new
head, using $O(1)$ extra space.}

\begin{patternbox}
Reversing a singly linked list is the archetypal \textbf{pointer-rewiring}
loop: walk the list carrying three pointers -- \texttt{prev},
\texttt{curr}, \texttt{next} -- and at each node flip \texttt{curr->next} to
point backward at \texttt{prev}.
\end{patternbox}

\subsection*{Flip each link, remember the rest}

The danger in rewiring is losing the remainder of the list. So before
flipping \texttt{curr->next} to \texttt{prev}, save \texttt{curr->next} in a
temporary \texttt{next}. Then set \texttt{curr->next = prev}, advance
\texttt{prev} to \texttt{curr}, and \texttt{curr} to the saved
\texttt{next}. When \texttt{curr} becomes null, \texttt{prev} is the new
head (the old tail).

\begin{ideabox}[Save Next Before You Overwrite It]
Each step overwrites \texttt{curr->next}, which is the only link to the rest
of the list -- so it must be stashed first. This save--flip--advance trio is
the canonical in-place reversal and recurs inside reverse-in-groups and
palindrome checks.
\end{ideabox}

\subsection*{Algorithm}

\begin{enumerate}
  \item \texttt{prev = null}, \texttt{curr = head}.
  \item While \texttt{curr}: save \texttt{next = curr->next}; set
        \texttt{curr->next = prev}; advance \texttt{prev = curr},
        \texttt{curr = next}.
  \item Return \texttt{prev}.
\end{enumerate}

\subsection*{Dry run}

\dryrun{List $1\to2\to3$.}

\begin{itemize}
  \item prev=null,curr=1: flip $1\to$null; prev=1,curr=2.
  \item flip $2\to1$; prev=2,curr=3. flip $3\to2$; prev=3,curr=null.
  \item New head $3\to2\to1$.
\end{itemize}

\subsection*{C++ implementation}

\begin{lstlisting}
#include <bits/stdc++.h>
using namespace std;

struct ListNode {
    int val; ListNode* next;
    ListNode(int v) : val(v), next(nullptr) {}
};

ListNode* build(const vector<int>& a) {
    ListNode dummy(0), *t = &dummy;
    for (int x : a) { t->next = new ListNode(x); t = t->next; }
    return dummy.next;
}

ListNode* reverse(ListNode* head) {
    ListNode* prev = nullptr;
    ListNode* curr = head;
    while (curr) {
        ListNode* next = curr->next;              // save the rest
        curr->next = prev;                        // flip
        prev = curr;                              // advance
        curr = next;
    }
    return prev;                                  // new head
}

void print(ListNode* h) { for (; h; h = h->next) cout << h->val << " "; cout << "\n"; }

int main() {
    ListNode* head = build({1, 2, 3, 4});
    head = reverse(head);
    print(head);                                  // 4 3 2 1
    return 0;
}
\end{lstlisting}

\begin{complexitybox}
\textbf{Time Complexity.} $O(n)$. \textbf{Space Complexity.} $O(1)$.
\end{complexitybox}

\begin{takeawaybox}
Iterative reversal is the save--flip--advance loop with three pointers,
$O(1)$ space. It is the most important pointer-manipulation pattern in the
chapter -- committed to muscle memory, it makes reverse-in-$k$-groups,
palindrome, and add-two-numbers straightforward.
\end{takeawaybox}

\section{Reverse a Linked List (Recursive)}
\label{sec:ll-reverse-recursive}

\problemstatement{Reverse a singly linked list using recursion, returning
the new head.}

\begin{patternbox}
The recursive reversal expresses the same rewiring differently: reverse the
rest of the list first, then fix the link between the current node and its
successor on the way back up. The new head bubbles up unchanged from the
deepest call.
\end{patternbox}

\subsection*{Reverse the tail, then flip one link}

Recurse to the end: the base case (empty or single node) returns itself as
the new head. On the way back, the current node's successor
(\texttt{head->next}) should now point back at the current node, so set
\texttt{head->next->next = head} and \texttt{head->next = null}. The deepest
call's return value -- the original last node -- is the new head and is
passed up unchanged through every frame.

\begin{ideabox}[Fix the Link on the Return Trip]
Going down, we only recurse; coming up, each frame flips exactly one link
between \texttt{head} and \texttt{head->next}. Because the new head is
determined at the very bottom, it is threaded back up untouched while the
per-frame flips reverse the pointers.
\end{ideabox}

\subsection*{Algorithm}

\begin{enumerate}
  \item If \texttt{head} is null or \texttt{head->next} is null, return
        \texttt{head}.
  \item \texttt{newHead = reverse(head->next)}.
  \item \texttt{head->next->next = head}; \texttt{head->next = null}.
  \item Return \texttt{newHead}.
\end{enumerate}

\subsection*{Dry run}

\dryrun{List $1\to2\to3$.}

\begin{itemize}
  \item Recurse to $3$ (base) $\to$ newHead $=3$.
  \item At $2$: $3\to2$, $2\to$null. At $1$: $2\to1$, $1\to$null.
  \item Result $3\to2\to1$.
\end{itemize}

\subsection*{C++ implementation}

\begin{lstlisting}
#include <bits/stdc++.h>
using namespace std;

struct ListNode {
    int val; ListNode* next;
    ListNode(int v) : val(v), next(nullptr) {}
};

ListNode* build(const vector<int>& a) {
    ListNode dummy(0), *t = &dummy;
    for (int x : a) { t->next = new ListNode(x); t = t->next; }
    return dummy.next;
}

ListNode* reverse(ListNode* head) {
    if (!head || !head->next) return head;        // base case
    ListNode* newHead = reverse(head->next);
    head->next->next = head;                      // flip one link
    head->next = nullptr;
    return newHead;                               // threaded up unchanged
}

void print(ListNode* h) { for (; h; h = h->next) cout << h->val << " "; cout << "\n"; }

int main() {
    ListNode* head = build({1, 2, 3, 4});
    head = reverse(head);
    print(head);                                  // 4 3 2 1
    return 0;
}
\end{lstlisting}

\begin{complexitybox}
\textbf{Time Complexity.} $O(n)$. \textbf{Space Complexity.} $O(n)$ for the
recursion stack (versus $O(1)$ for the iterative form).
\end{complexitybox}

\begin{takeawaybox}
Recursive reversal reverses the tail then flips one link per return frame,
carrying the new head up from the base case. It is more elegant but uses
$O(n)$ stack; prefer the iterative version when depth could overflow.
\end{takeawaybox}

\section{Detect a Loop in a Linked List}
\label{sec:ll-detect-loop}

\problemstatement{Given the head of a singly linked list, determine whether
it contains a cycle (some node's \texttt{next} points back to an earlier
node), using $O(1)$ extra space.}

\begin{patternbox}
\emph{``Is there a cycle?''} is \textbf{Floyd's tortoise--hare}: a slow
pointer (one step) and a fast pointer (two steps). If there is a loop, the
fast pointer laps the slow one and they meet inside the loop; if there is no
loop, the fast pointer reaches null.
\end{patternbox}

\subsection*{Why the two pointers must meet}

Inside a loop, each iteration the fast pointer gains one node on the slow
pointer, so the gap shrinks by one every step and eventually hits zero --
they collide. If the list is acyclic, the fast pointer falls off the end
(reaches null) first. This detects a cycle in one pass with only two
pointers, unlike a hash-set approach that needs $O(n)$ memory.

\begin{ideabox}[Gaining One Node per Step Guarantees a Meeting]
Relative to the slow pointer, the fast pointer moves one node closer each
iteration. On a finite cycle the distance between them is bounded, so it
must reach zero -- a collision -- proving a loop exists.
\end{ideabox}

\subsection*{Algorithm}

\begin{enumerate}
  \item \texttt{slow = fast = head}.
  \item While \texttt{fast} and \texttt{fast->next}: advance \texttt{slow}
        by one, \texttt{fast} by two; if they are equal, return true.
  \item If the loop exits, return false.
\end{enumerate}

\subsection*{Dry run}

\dryrun{$1\to2\to3\to4\to2$ (node $4$ links back to $2$).}

\begin{itemize}
  \item slow/fast advance; fast laps and meets slow at some node inside
        $\{2,3,4\}$.
  \item Meeting confirms a cycle -- return true.
\end{itemize}

\subsection*{C++ implementation}

\begin{lstlisting}
#include <bits/stdc++.h>
using namespace std;

struct ListNode {
    int val; ListNode* next;
    ListNode(int v) : val(v), next(nullptr) {}
};

bool hasCycle(ListNode* head) {
    ListNode* slow = head;
    ListNode* fast = head;
    while (fast && fast->next) {
        slow = slow->next;
        fast = fast->next->next;
        if (slow == fast) return true;            // collision inside loop
    }
    return false;
}

int main() {
    ListNode* a = new ListNode(1);
    a->next = new ListNode(2);
    a->next->next = new ListNode(3);
    a->next->next->next = new ListNode(4);
    a->next->next->next->next = a->next;          // 4 -> 2, a cycle
    cout << (hasCycle(a) ? "cycle" : "no cycle") << "\n";
    return 0;
}
\end{lstlisting}

\begin{complexitybox}
\textbf{Time Complexity.} $O(n)$. \textbf{Space Complexity.} $O(1)$ (a hash
set would cost $O(n)$).
\end{complexitybox}

\begin{takeawaybox}
Floyd's cycle detection uses two speeds and constant memory; a collision
proves a loop. The same meeting point locates the loop's start and length
(next two sections), so this is the gateway to all cycle problems.
\end{takeawaybox}

\section{Starting Point of a Loop}
\label{sec:ll-loop-start}

\problemstatement{If a linked list has a cycle, return the node where the
cycle begins; otherwise return null. Use $O(1)$ extra space.}

\begin{patternbox}
After Floyd's slow/fast pointers meet inside the loop, a beautiful fact
kicks in: \textbf{resetting one pointer to the head and advancing both one
step at a time, they meet exactly at the loop's start}. It follows from the
distance arithmetic of the tortoise--hare.
\end{patternbox}

\subsection*{The distance argument}

Let the distance from head to the loop start be $L$, and let the slow and
fast pointers meet $k$ nodes into the loop. When they meet, the fast
pointer has travelled twice the slow pointer's distance; working through the
algebra shows that $L$ equals the remaining distance from the meeting point
around to the loop start. So if you place one pointer at the head and leave
the other at the meeting point, and move both one step at a time, they
converge precisely at the loop's entrance.

\begin{ideabox}[Head to Start Equals Meeting Point to Start]
The tortoise--hare meeting point is positioned so that the number of steps
from the head to the loop start equals the number of steps from the meeting
point (going forward around the loop) to the loop start. Two synchronized
one-step walks therefore collide at the start.
\end{ideabox}

\subsection*{Algorithm}

\begin{enumerate}
  \item Run slow/fast until they meet; if fast reaches null, return null (no
        loop).
  \item Move \texttt{slow} to the head. Advance \texttt{slow} and
        \texttt{fast} one step each until they are equal.
  \item That node is the loop's start.
\end{enumerate}

\subsection*{Dry run}

\dryrun{$1\to2\to3\to4\to2$ (loop starts at $2$).}

\begin{itemize}
  \item Slow/fast meet inside the loop; reset slow to head ($1$).
  \item One step each: they converge at node $2$ -- the loop start.
\end{itemize}

\subsection*{C++ implementation}

\begin{lstlisting}
#include <bits/stdc++.h>
using namespace std;

struct ListNode {
    int val; ListNode* next;
    ListNode(int v) : val(v), next(nullptr) {}
};

ListNode* detectCycle(ListNode* head) {
    ListNode* slow = head;
    ListNode* fast = head;
    while (fast && fast->next) {
        slow = slow->next;
        fast = fast->next->next;
        if (slow == fast) {                       // meeting point
            slow = head;
            while (slow != fast) { slow = slow->next; fast = fast->next; }
            return slow;                          // loop start
        }
    }
    return nullptr;
}

int main() {
    ListNode* a = new ListNode(1);
    a->next = new ListNode(2);
    a->next->next = new ListNode(3);
    a->next->next->next = new ListNode(4);
    a->next->next->next->next = a->next;          // 4 -> 2
    ListNode* start = detectCycle(a);
    cout << (start ? start->val : -1) << "\n";     // 2
    return 0;
}
\end{lstlisting}

\begin{complexitybox}
\textbf{Time Complexity.} $O(n)$. \textbf{Space Complexity.} $O(1)$.
\end{complexitybox}

\begin{takeawaybox}
The loop's start is found by resetting one pointer to the head after the
tortoise--hare collision and walking both one step at a time. The
underlying ``$L$ equals meeting-to-start'' identity is worth remembering --
it is asked as a follow-up to cycle detection constantly.
\end{takeawaybox}

\section{Length of a Loop}
\label{sec:ll-loop-length}

\problemstatement{If a linked list contains a cycle, return the number of
nodes in the cycle; otherwise return $0$.}

\begin{patternbox}
Once Floyd's pointers meet inside the loop, \textbf{keep one pointer fixed
at the meeting node and walk the other around until it returns} -- the
number of steps is the loop length. The collision guarantees both pointers
are on the cycle.
\end{patternbox}

\subsection*{Walk once around from the meeting point}

Run the tortoise--hare until \texttt{slow} and \texttt{fast} coincide; that
node lies on the cycle. Now advance a single pointer from that node,
counting steps, until it comes back to the same node. Since it started on
the cycle, it traverses exactly the cycle once, and the count is the loop's
length. If the pointers never meet (fast reaches null), the length is $0$.

\begin{ideabox}[A Full Lap from Any Cycle Node Measures It]
Any node known to be on the cycle can measure the cycle's size: leave it in
place and count steps until a moving pointer returns to it. The Floyd
meeting point conveniently provides such a node.
\end{ideabox}

\subsection*{Algorithm}

\begin{enumerate}
  \item Run slow/fast until they meet; if no meeting, return $0$.
  \item From the meeting node, advance a pointer counting steps until it
        returns to the meeting node.
  \item Return the count.
\end{enumerate}

\subsection*{Dry run}

\dryrun{$1\to2\to3\to4\to2$ (loop $2\to3\to4\to2$, length $3$).}

\begin{itemize}
  \item Pointers meet inside $\{2,3,4\}$. From the meeting node, count a
        lap: $3$ steps back to it.
  \item Length $=3$.
\end{itemize}

\subsection*{C++ implementation}

\begin{lstlisting}
#include <bits/stdc++.h>
using namespace std;

struct ListNode {
    int val; ListNode* next;
    ListNode(int v) : val(v), next(nullptr) {}
};

int loopLength(ListNode* head) {
    ListNode* slow = head;
    ListNode* fast = head;
    while (fast && fast->next) {
        slow = slow->next;
        fast = fast->next->next;
        if (slow == fast) {                       // on the cycle
            int len = 1;
            ListNode* p = slow->next;
            while (p != slow) { p = p->next; len++; }
            return len;
        }
    }
    return 0;
}

int main() {
    ListNode* a = new ListNode(1);
    a->next = new ListNode(2);
    a->next->next = new ListNode(3);
    a->next->next->next = new ListNode(4);
    a->next->next->next->next = a->next;          // 4 -> 2
    cout << loopLength(a) << "\n";                 // 3
    return 0;
}
\end{lstlisting}

\begin{complexitybox}
\textbf{Time Complexity.} $O(n)$. \textbf{Space Complexity.} $O(1)$.
\end{complexitybox}

\begin{takeawaybox}
Loop length is one lap counted from the tortoise--hare meeting node. Cycle
detection, its start, and its length are three questions answered by the
same slow/fast collision -- keep the trio together in memory.
\end{takeawaybox}

\section{Check if a Linked List is a Palindrome}
\label{sec:ll-palindrome}

\problemstatement{Given the head of a singly linked list, determine whether
the sequence of values reads the same forwards and backwards, using $O(1)$
extra space.}

\begin{patternbox}
Without random access, comparing front to back means \textbf{finding the
middle (slow/fast), reversing the second half, then walking both halves in
step}. Equality of the two halves proves a palindrome. This combines two
earlier techniques into one.
\end{patternbox}

\subsection*{Middle, reverse, compare}

Use the tortoise--hare to reach the middle. Reverse the second half in place
(the iterative reversal). Now walk one pointer from the head and another
from the reversed second half, comparing values; if all match, the list is a
palindrome. Optionally restore the list by reversing the second half back.
This achieves $O(1)$ space, versus copying values into an array ($O(n)$).

\begin{ideabox}[Reverse Half to Compare Without Extra Memory]
A palindrome check needs to read the list backward, which a singly linked
list cannot do -- so reverse the back half physically. Comparing the
untouched front half against the reversed back half is exactly the
front-to-back comparison, in constant space.
\end{ideabox}

\subsection*{Algorithm}

\begin{enumerate}
  \item Find the middle with slow/fast.
  \item Reverse the second half.
  \item Compare the first half and the reversed second half node by node.
  \item Return whether all comparisons matched.
\end{enumerate}

\subsection*{Dry run}

\dryrun{List $1\to2\to2\to1$.}

\begin{itemize}
  \item Middle after node $2$ (second); reverse back half $\to1\to2$.
  \item Compare $1{=}1$, $2{=}2$ -- palindrome.
\end{itemize}

\subsection*{C++ implementation}

\begin{lstlisting}
#include <bits/stdc++.h>
using namespace std;

struct ListNode {
    int val; ListNode* next;
    ListNode(int v) : val(v), next(nullptr) {}
};

ListNode* build(const vector<int>& a) {
    ListNode dummy(0), *t = &dummy;
    for (int x : a) { t->next = new ListNode(x); t = t->next; }
    return dummy.next;
}

ListNode* reverse(ListNode* head) {
    ListNode* prev = nullptr;
    while (head) { ListNode* nx = head->next; head->next = prev; prev = head; head = nx; }
    return prev;
}

bool isPalindrome(ListNode* head) {
    ListNode* slow = head;
    ListNode* fast = head;
    while (fast && fast->next) { slow = slow->next; fast = fast->next->next; }
    ListNode* second = reverse(slow);             // reverse back half

    ListNode* p1 = head;
    ListNode* p2 = second;
    while (p2) {
        if (p1->val != p2->val) return false;
        p1 = p1->next; p2 = p2->next;
    }
    return true;
}

int main() {
    cout << isPalindrome(build({1, 2, 2, 1})) << "\n";   // 1
    cout << isPalindrome(build({1, 2, 3})) << "\n";      // 0
    return 0;
}
\end{lstlisting}

\begin{complexitybox}
\textbf{Time Complexity.} $O(n)$ -- find middle, reverse, compare, each
linear. \textbf{Space Complexity.} $O(1)$.
\end{complexitybox}

\begin{takeawaybox}
Palindrome checking composes slow/fast (middle) with in-place reversal
(back half) to compare front and back in constant space. It is the classic
demonstration that combining two basic list techniques beats the $O(n)$-
memory array copy.
\end{takeawaybox}

\section{Segregate Odd and Even Position Nodes}
\label{sec:ll-odd-even}

\problemstatement{Given a linked list, group all nodes at odd positions
together followed by the nodes at even positions, preserving the relative
order within each group. Positions are 1-indexed; do it in place.}

\begin{patternbox}
Rearranging by position parity is a \textbf{two-chain weave}: maintain an
odd chain and an even chain, appending alternate nodes to each, then splice
the even chain onto the tail of the odd chain. Only pointers move.
\end{patternbox}

\subsection*{Build two chains, then join}

Keep an \texttt{odd} pointer starting at the head and an \texttt{even}
pointer starting at the second node, plus a saved \texttt{evenHead}. Repeatedly
link \texttt{odd->next = even->next} (skip the even node) and advance
\texttt{odd}; then \texttt{even->next = odd->next} and advance
\texttt{even}. This splits the list into odd-position and even-position
chains in a single pass. Finally attach \texttt{evenHead} after the last odd
node.

\begin{ideabox}[Interleave by Skipping Every Other Node]
Advancing each chain by two positions (skipping the other parity) threads
the odd and even nodes apart without moving data. The saved
\texttt{evenHead} lets the two chains be reconnected at the end.
\end{ideabox}

\subsection*{Algorithm}

\begin{enumerate}
  \item If the list has fewer than $3$ nodes, return it unchanged.
  \item \texttt{odd = head}, \texttt{even = head->next},
        \texttt{evenHead = even}.
  \item While \texttt{even} and \texttt{even->next}: relink odd then even by
        skipping one node each; advance both.
  \item \texttt{odd->next = evenHead}; return the head.
\end{enumerate}

\subsection*{Dry run}

\dryrun{$1\to2\to3\to4\to5$.}

\begin{itemize}
  \item Odd chain $1\to3\to5$; even chain $2\to4$.
  \item Join: $1\to3\to5\to2\to4$.
\end{itemize}

\subsection*{C++ implementation}

\begin{lstlisting}
#include <bits/stdc++.h>
using namespace std;

struct ListNode {
    int val; ListNode* next;
    ListNode(int v) : val(v), next(nullptr) {}
};

ListNode* build(const vector<int>& a) {
    ListNode dummy(0), *t = &dummy;
    for (int x : a) { t->next = new ListNode(x); t = t->next; }
    return dummy.next;
}

ListNode* oddEvenList(ListNode* head) {
    if (!head || !head->next) return head;
    ListNode* odd = head;
    ListNode* even = head->next;
    ListNode* evenHead = even;
    while (even && even->next) {
        odd->next = even->next;   odd = odd->next;    // splice next odd
        even->next = odd->next;   even = even->next;  // splice next even
    }
    odd->next = evenHead;                            // attach evens after odds
    return head;
}

void print(ListNode* h) { for (; h; h = h->next) cout << h->val << " "; cout << "\n"; }

int main() {
    print(oddEvenList(build({1, 2, 3, 4, 5})));    // 1 3 5 2 4
    return 0;
}
\end{lstlisting}

\begin{complexitybox}
\textbf{Time Complexity.} $O(n)$ -- one pass. \textbf{Space Complexity.}
$O(1)$.
\end{complexitybox}

\begin{takeawaybox}
Grouping by position parity is a two-chain weave joined at the end -- a
pure pointer rearrangement. The ``split into interleaved chains, then
reconnect'' idea also underlies reordering and unzip-style list problems.
\end{takeawaybox}

\section{Remove the Nth Node from the End}
\label{sec:ll-remove-nth}

\problemstatement{Given the head of a linked list and an integer $n$, remove
the $n$-th node counting from the end and return the head, in one pass.}

\begin{patternbox}
\emph{``$n$-th from the end in one pass''} is a \textbf{two-pointer gap}:
advance a lead pointer $n$ nodes ahead, then move both until the lead
reaches the end -- the trailing pointer now sits just before the node to
delete. A dummy head handles removing the first node cleanly.
\end{patternbox}

\subsection*{Fix the gap, then slide it}

Place a dummy node before the head. Advance a \texttt{fast} pointer $n$ steps
from the dummy. Then move \texttt{fast} and a \texttt{slow} pointer (starting
at the dummy) together until \texttt{fast} reaches the last node. Because the
gap between them is exactly $n$, \texttt{slow} ends up just before the target;
unlink it with \texttt{slow->next = slow->next->next}. Return
\texttt{dummy.next}.

\begin{ideabox}[A Fixed Gap Finds the Position from the End]
The distance from the end is converted into a fixed lead of $n$ nodes.
Sliding both pointers until the lead hits the tail leaves the trailing
pointer at the predecessor of the $n$-th-from-last node -- all in a single
traversal.
\end{ideabox}

\subsection*{Algorithm}

\begin{enumerate}
  \item Dummy before head; \texttt{fast = dummy}, advance it $n$ steps.
  \item Move \texttt{fast} and \texttt{slow} (from dummy) together until
        \texttt{fast->next} is null.
  \item \texttt{slow->next = slow->next->next}; return \texttt{dummy.next}.
\end{enumerate}

\subsection*{Dry run}

\dryrun{$1\to2\to3\to4\to5$, $n=2$.}

\begin{itemize}
  \item fast advances $2$ to node $2$; then both slide until fast at $5$.
  \item slow at $3$; remove $3->next$ (node $4$). Result $1\to2\to3\to5$.
\end{itemize}

\subsection*{C++ implementation}

\begin{lstlisting}
#include <bits/stdc++.h>
using namespace std;

struct ListNode {
    int val; ListNode* next;
    ListNode(int v) : val(v), next(nullptr) {}
};

ListNode* build(const vector<int>& a) {
    ListNode dummy(0), *t = &dummy;
    for (int x : a) { t->next = new ListNode(x); t = t->next; }
    return dummy.next;
}

ListNode* removeNthFromEnd(ListNode* head, int n) {
    ListNode dummy(0); dummy.next = head;
    ListNode* fast = &dummy;
    for (int i = 0; i < n; i++) fast = fast->next;   // lead by n
    ListNode* slow = &dummy;
    while (fast->next) { fast = fast->next; slow = slow->next; }
    slow->next = slow->next->next;                   // unlink target
    return dummy.next;
}

void print(ListNode* h) { for (; h; h = h->next) cout << h->val << " "; cout << "\n"; }

int main() {
    print(removeNthFromEnd(build({1, 2, 3, 4, 5}), 2));   // 1 2 3 5
    return 0;
}
\end{lstlisting}

\begin{complexitybox}
\textbf{Time Complexity.} $O(n)$ -- a single pass. \textbf{Space
Complexity.} $O(1)$.
\end{complexitybox}

\begin{takeawaybox}
A fixed two-pointer gap of $n$ turns ``from the end'' into a one-pass
position. The dummy node makes deleting the head uniform with any other
node -- the recurring reason to prepend a dummy.
\end{takeawaybox}

\section{Delete the Middle Node}
\label{sec:ll-delete-middle}

\problemstatement{Given the head of a linked list, delete the middle node
(for even length, the second middle) and return the head. If the list has
one node, return null.}

\begin{patternbox}
Deleting the middle combines the slow/fast middle-finder with the need for
the node \emph{before} the middle. Advance \texttt{fast} by two but start it
one step ahead (or track \texttt{prev}) so that when \texttt{fast} reaches
the end, \texttt{slow} stops just before the middle, ready to unlink it.
\end{patternbox}

\subsection*{Stop slow one node early}

Run the tortoise--hare, but arrange the offsets so \texttt{slow} lands on the
\emph{predecessor} of the middle rather than the middle itself -- for
instance by giving \texttt{fast} a two-node head start. Then
\texttt{slow->next = slow->next->next} removes the middle node. Handle the
single-node list as a special case (return null).

\begin{ideabox}[Delete Needs the Predecessor]
Finding the middle is not enough to delete it in a singly linked list; you
need its predecessor to rewire \texttt{next}. Offsetting the fast pointer so
slow halts one node early yields that predecessor directly, avoiding a
second pass.
\end{ideabox}

\subsection*{Algorithm}

\begin{enumerate}
  \item If the list has one node, return null.
  \item \texttt{slow = head}, \texttt{fast = head->next->next}.
  \item While \texttt{fast} and \texttt{fast->next}: advance \texttt{slow}
        by one, \texttt{fast} by two.
  \item \texttt{slow->next = slow->next->next}; return the head.
\end{enumerate}

\subsection*{Dry run}

\dryrun{$1\to2\to3\to4$ (delete second middle, node $3$).}

\begin{itemize}
  \item slow$=1$, fast$=3$. Step: slow$=2$, fast$=$null.
  \item Remove $2->next$ (node $3$). Result $1\to2\to4$.
\end{itemize}

\subsection*{C++ implementation}

\begin{lstlisting}
#include <bits/stdc++.h>
using namespace std;

struct ListNode {
    int val; ListNode* next;
    ListNode(int v) : val(v), next(nullptr) {}
};

ListNode* build(const vector<int>& a) {
    ListNode dummy(0), *t = &dummy;
    for (int x : a) { t->next = new ListNode(x); t = t->next; }
    return dummy.next;
}

ListNode* deleteMiddle(ListNode* head) {
    if (!head->next) return nullptr;              // single node
    ListNode* slow = head;
    ListNode* fast = head->next->next;            // two-node head start
    while (fast && fast->next) {
        slow = slow->next;
        fast = fast->next->next;
    }
    slow->next = slow->next->next;                // unlink the middle
    return head;
}

void print(ListNode* h) { for (; h; h = h->next) cout << h->val << " "; cout << "\n"; }

int main() {
    print(deleteMiddle(build({1, 2, 3, 4, 5})));   // 1 2 4 5
    print(deleteMiddle(build({1, 2, 3, 4})));      // 1 2 4
    return 0;
}
\end{lstlisting}

\begin{complexitybox}
\textbf{Time Complexity.} $O(n)$. \textbf{Space Complexity.} $O(1)$.
\end{complexitybox}

\begin{takeawaybox}
To delete the middle, offset the fast pointer so slow halts at the middle's
predecessor -- then a single rewire removes it. Adjusting the two-pointer
offset to land on the predecessor is a common tweak to the middle-finder.
\end{takeawaybox}

\section{Sort a Linked List}
\label{sec:ll-sort-list}

\problemstatement{Given the head of a linked list, sort it in ascending
order and return the new head, ideally in $O(n\log n)$ time and $O(1)$
auxiliary space.}

\begin{patternbox}
\textbf{Merge sort} is the natural sort for linked lists: splitting by the
middle (slow/fast) and merging two sorted lists are both cheap pointer
operations, and unlike arrays no random access is needed. Quicksort's
partition is awkward on lists, so merge sort is the standard choice.
\end{patternbox}

\subsection*{Split, sort halves, merge}

Find the middle with the tortoise--hare and cut the list into two halves.
Recursively sort each half. Then merge the two sorted halves by repeatedly
appending the smaller current node -- exactly the merge step of array merge
sort, but relinking nodes instead of copying values. The recursion bottoms
out at single-node (already sorted) lists.

\begin{ideabox}[Lists Merge and Split Without Extra Arrays]
Splitting a list is just severing it at the middle; merging is relinking
nodes in order. Both are $O(1)$-space pointer work, so merge sort on a list
needs no auxiliary array -- its recursion stack is the only extra space.
\end{ideabox}

\subsection*{Algorithm}

\begin{enumerate}
  \item If the list has $0$ or $1$ nodes, return it.
  \item Split at the middle into two halves.
  \item Recursively sort both halves; merge them into one sorted list.
\end{enumerate}

\subsection*{Dry run}

\dryrun{$3\to1\to2$.}

\begin{itemize}
  \item Split into $3\to1$ and $2$. Sort $3\to1$ into $1\to3$.
  \item Merge $1\to3$ with $2$: $1\to2\to3$.
\end{itemize}

\subsection*{C++ implementation}

\begin{lstlisting}
#include <bits/stdc++.h>
using namespace std;

struct ListNode {
    int val; ListNode* next;
    ListNode(int v) : val(v), next(nullptr) {}
};

ListNode* build(const vector<int>& a) {
    ListNode dummy(0), *t = &dummy;
    for (int x : a) { t->next = new ListNode(x); t = t->next; }
    return dummy.next;
}

ListNode* merge(ListNode* a, ListNode* b) {
    ListNode dummy(0), *t = &dummy;
    while (a && b) {
        if (a->val <= b->val) { t->next = a; a = a->next; }
        else                  { t->next = b; b = b->next; }
        t = t->next;
    }
    t->next = a ? a : b;                          // attach the remainder
    return dummy.next;
}

ListNode* sortList(ListNode* head) {
    if (!head || !head->next) return head;
    ListNode* slow = head; ListNode* fast = head->next;
    while (fast && fast->next) { slow = slow->next; fast = fast->next->next; }
    ListNode* mid = slow->next;
    slow->next = nullptr;                         // cut into two halves
    return merge(sortList(head), sortList(mid));
}

void print(ListNode* h) { for (; h; h = h->next) cout << h->val << " "; cout << "\n"; }

int main() {
    print(sortList(build({4, 2, 1, 3})));         // 1 2 3 4
    return 0;
}
\end{lstlisting}

\begin{complexitybox}
\textbf{Time Complexity.} $O(n\log n)$ -- $\log n$ split levels, $O(n)$
merge per level. \textbf{Space Complexity.} $O(\log n)$ recursion stack.
\end{complexitybox}

\begin{takeawaybox}
Merge sort suits linked lists because split and merge are pointer
operations needing no random access or extra array. It is the default
answer to ``sort a list'' and reuses the middle-finder and merge routines
seen elsewhere.
\end{takeawaybox}

\section{Sort a Linked List of 0s, 1s, and 2s}
\label{sec:ll-sort-012}

\problemstatement{Given a linked list whose values are only $0$, $1$, or
$2$, sort it in a single pass without swapping data, returning the head.}

\begin{patternbox}
With just three distinct values, sorting is a \textbf{three-chain
partition}: build separate chains for $0$s, $1$s, and $2$s in one pass, then
concatenate them. No comparisons or $O(n\log n)$ sort are needed -- this is
the linked-list Dutch-national-flag.
\end{patternbox}

\subsection*{Three dummy heads, then join}

Create three dummy-headed chains, one per value. Walk the list once,
appending each node to the chain matching its value. Then link the $0$-chain
to the $1$-chain to the $2$-chain, terminating the last with null. Because
you distribute nodes by value and reassemble in order, the result is sorted.
Using dummy heads avoids special-casing empty chains.

\begin{ideabox}[Bucket by Value, Then Concatenate]
When the number of distinct keys is a small constant, a counting/bucketing
pass beats comparison sorting. Three chains bucket the nodes; concatenation
orders them -- linear time, no data movement, just relinking.
\end{ideabox}

\subsection*{Algorithm}

\begin{enumerate}
  \item Three dummy-headed chains for $0$, $1$, $2$; keep their tails.
  \item Walk the list, appending each node to its value's chain.
  \item Join $0\to1\to2$ (skipping empty chains via the dummy links); null-
        terminate; return the combined head.
\end{enumerate}

\subsection*{Dry run}

\dryrun{$1\to2\to0\to1\to2\to0$.}

\begin{itemize}
  \item Chains: $0$s $=0\to0$, $1$s $=1\to1$, $2$s $=2\to2$.
  \item Join: $0\to0\to1\to1\to2\to2$.
\end{itemize}

\subsection*{C++ implementation}

\begin{lstlisting}
#include <bits/stdc++.h>
using namespace std;

struct ListNode {
    int val; ListNode* next;
    ListNode(int v) : val(v), next(nullptr) {}
};

ListNode* build(const vector<int>& a) {
    ListNode dummy(0), *t = &dummy;
    for (int x : a) { t->next = new ListNode(x); t = t->next; }
    return dummy.next;
}

ListNode* sortColors(ListNode* head) {
    ListNode d0(0), d1(0), d2(0);                 // dummy heads
    ListNode *t0 = &d0, *t1 = &d1, *t2 = &d2;
    for (ListNode* c = head; c; c = c->next) {
        if (c->val == 0)      { t0->next = c; t0 = c; }
        else if (c->val == 1) { t1->next = c; t1 = c; }
        else                  { t2->next = c; t2 = c; }
    }
    t2->next = nullptr;                           // terminate
    t1->next = d2.next;                           // join 1s -> 2s
    t0->next = d1.next ? d1.next : d2.next;       // join 0s -> (1s or 2s)
    return d0.next;
}

void print(ListNode* h) { for (; h; h = h->next) cout << h->val << " "; cout << "\n"; }

int main() {
    print(sortColors(build({1, 2, 0, 1, 2, 0})));  // 0 0 1 1 2 2
    return 0;
}
\end{lstlisting}

\begin{complexitybox}
\textbf{Time Complexity.} $O(n)$ -- one pass. \textbf{Space Complexity.}
$O(1)$ (three dummy nodes).
\end{complexitybox}

\begin{takeawaybox}
A constant number of distinct values makes bucketing beat comparison
sorting: three chains, one pass, then concatenate. The dummy-headed chain
pattern -- collect into groups, then splice -- reappears in partitioning and
stable-reordering problems.
\end{takeawaybox}

\section{Intersection Point of Two Linked Lists}
\label{sec:ll-intersection}

\problemstatement{Given the heads of two singly linked lists that may merge
at some node (forming a Y shape), return the first common node, or null if
they never intersect. Use $O(1)$ extra space.}

\begin{patternbox}
The elegant solution is the \textbf{two-pointer switch}: walk one pointer
along list A then list B, and another along list B then list A. They travel
the same total distance, so they arrive at the intersection simultaneously
(or at null together).
\end{patternbox}

\subsection*{Equalise the path lengths automatically}

Let the lists have lengths $a$ and $b$ with a shared tail of length $c$.
Pointer 1 walks $a$ then $b$; pointer 2 walks $b$ then $a$; both cover
$a+b$ nodes. After the switch, each is exactly $c$ nodes from the end, so
they meet at the first shared node. If there is no intersection, both reach
null at the same time. This removes the need to measure and align lengths
manually.

\begin{ideabox}[Switching Lists Cancels the Length Difference]
By having each pointer traverse both lists, the head-start difference $|a-b|$
is absorbed: after the switch they are equidistant from the merge point. The
first place they coincide is the intersection.
\end{ideabox}

\subsection*{Algorithm}

\begin{enumerate}
  \item \texttt{p = headA}, \texttt{q = headB}.
  \item Advance both; when \texttt{p} hits null, redirect it to
        \texttt{headB}; when \texttt{q} hits null, redirect to
        \texttt{headA}.
  \item Stop when \texttt{p == q}; that node (or null) is the answer.
\end{enumerate}

\subsection*{Dry run}

\dryrun{A $=1\to9\to7$, B $=6\to7$, sharing the node $7$.}

\begin{itemize}
  \item p: $1,9,7,\text{(B)}6,7$. q: $6,7,\text{(A)}1,9,7$.
  \item Both reach the shared $7$ at the same step -- the intersection.
\end{itemize}

\subsection*{C++ implementation}

\begin{lstlisting}
#include <bits/stdc++.h>
using namespace std;

struct ListNode {
    int val; ListNode* next;
    ListNode(int v) : val(v), next(nullptr) {}
};

ListNode* getIntersectionNode(ListNode* a, ListNode* b) {
    if (!a || !b) return nullptr;
    ListNode* p = a; ListNode* q = b;
    while (p != q) {
        p = p ? p->next : b;                      // switch to B at the end
        q = q ? q->next : a;                      // switch to A at the end
    }
    return p;                                     // meeting node or null
}

int main() {
    ListNode* shared = new ListNode(7);
    shared->next = new ListNode(8);
    ListNode* a = new ListNode(1);
    a->next = new ListNode(9); a->next->next = shared;
    ListNode* b = new ListNode(6); b->next = shared;

    ListNode* r = getIntersectionNode(a, b);
    cout << (r ? r->val : -1) << "\n";             // 7
    return 0;
}
\end{lstlisting}

\begin{complexitybox}
\textbf{Time Complexity.} $O(a+b)$ -- each pointer walks both lists once.
\textbf{Space Complexity.} $O(1)$ (a hash set would cost $O(n)$).
\end{complexitybox}

\begin{takeawaybox}
The two-pointer list switch equalises unequal lengths automatically and
finds the intersection in constant space. It is a favourite interview
trick precisely because the length-cancellation is so clean.
\end{takeawaybox}

\section{Add One to a Number Represented by a Linked List}
\label{sec:ll-add-one}

\problemstatement{A non-negative integer is stored as a linked list, most
significant digit first (head). Add one to the number and return the head of
the resulting list.}

\begin{patternbox}
The carry propagates from the \emph{least} significant digit, but the list
is stored most-significant first -- so either \textbf{reverse, add, reverse
back}, or use recursion to reach the last node and carry upward on the
return trip. Both avoid needing backward pointers.
\end{patternbox}

\subsection*{Carry from the back via recursion}

Recurse to the tail; the deepest call adds one and returns any carry. Each
frame adds the incoming carry to its digit, updates the digit modulo $10$,
and returns the new carry ($1$ if the digit was $9$, else $0$). If the head
still produces a carry after processing, prepend a new leading $1$ node.
This treats the recursion stack as the ``reversed'' traversal, so no
explicit reversal is required.

\begin{ideabox}[The Return Trip Is the Least-Significant-First Pass]
Recursion reaches the last (least significant) digit first on the way down,
then processes digits back-to-front on the way up -- exactly the order
addition needs. A leftover carry at the very top means a new most-
significant digit.
\end{ideabox}

\subsection*{Algorithm}

\begin{enumerate}
  \item Recurse to the last node; add $1$ there, return the carry.
  \item Each frame: \texttt{sum = digit + carry}; set digit to
        \texttt{sum \% 10}; return \texttt{sum / 10}.
  \item If the top-level carry is $1$, prepend a new node with value $1$.
\end{enumerate}

\subsection*{Dry run}

\dryrun{List $1\to2\to9$ (number $129$).}

\begin{itemize}
  \item Last digit $9+1=10$: digit $0$, carry $1$. Next $2+1=3$, carry $0$.
  \item Head $1$ unchanged. Result $1\to3\to0$ ($130$).
\end{itemize}

\subsection*{C++ implementation}

\begin{lstlisting}
#include <bits/stdc++.h>
using namespace std;

struct ListNode {
    int val; ListNode* next;
    ListNode(int v) : val(v), next(nullptr) {}
};

ListNode* build(const vector<int>& a) {
    ListNode dummy(0), *t = &dummy;
    for (int x : a) { t->next = new ListNode(x); t = t->next; }
    return dummy.next;
}

int addHelper(ListNode* node) {
    if (!node) return 1;                          // seed the +1 at the end
    int carry = addHelper(node->next);
    int sum = node->val + carry;
    node->val = sum % 10;
    return sum / 10;
}

ListNode* addOne(ListNode* head) {
    int carry = addHelper(head);
    if (carry) {                                  // new leading digit
        ListNode* node = new ListNode(carry);
        node->next = head;
        return node;
    }
    return head;
}

void print(ListNode* h) { for (; h; h = h->next) cout << h->val; cout << "\n"; }

int main() {
    print(addOne(build({1, 2, 9})));              // 130
    print(addOne(build({9, 9, 9})));              // 1000
    return 0;
}
\end{lstlisting}

\begin{complexitybox}
\textbf{Time Complexity.} $O(n)$. \textbf{Space Complexity.} $O(n)$
recursion stack (or $O(1)$ with the reverse--add--reverse approach).
\end{complexitybox}

\begin{takeawaybox}
Arithmetic on a most-significant-first list needs a back-to-front pass,
supplied for free by recursion's return trip (or by reversing). A leftover
carry at the head becomes a new leading digit -- the one edge case to
remember.
\end{takeawaybox}

\section{Add Two Numbers Represented by Linked Lists}
\label{sec:ll-add-two-numbers}

\problemstatement{Two non-negative integers are stored as linked lists with
the \emph{least} significant digit first. Add them and return the sum as a
linked list in the same order.}

\begin{patternbox}
Because the digits are stored least-significant first, addition proceeds
head-to-tail with a running \textbf{carry} -- a straightforward simultaneous
walk of both lists, building the result node by node, exactly like grade-
school addition.
\end{patternbox}

\subsection*{Walk both lists with a carry}

Traverse both lists together. At each step add the two current digits (or
$0$ if a list is exhausted) plus the carry; the new digit is the sum modulo
$10$ and the new carry is the sum divided by $10$. Append the digit to the
result via a dummy-headed builder. Continue until both lists end and the
carry is zero -- a final carry produces one more leading node.

\begin{ideabox}[Least-Significant-First Makes Addition a Forward Walk]
Storing digits reversed means the carry flows in the same direction as
traversal, so no reversal or recursion is needed -- just a single forward
pass summing aligned digits with a carry, the most natural form of the
problem.
\end{ideabox}

\subsection*{Algorithm}

\begin{enumerate}
  \item Dummy-headed result; \texttt{carry = 0}.
  \item While either list has nodes or \texttt{carry} is nonzero: sum the
        available digits and carry; append \texttt{sum \% 10}; set
        \texttt{carry = sum / 10}.
  \item Return \texttt{dummy.next}.
\end{enumerate}

\subsection*{Dry run}

\dryrun{$2\to4\to3$ ($342$) and $5\to6\to4$ ($465$).}

\begin{itemize}
  \item $2{+}5{=}7$; $4{+}6{=}10\to0$ carry $1$; $3{+}4{+}1{=}8$.
  \item Result $7\to0\to8$ ($807$).
\end{itemize}

\subsection*{C++ implementation}

\begin{lstlisting}
#include <bits/stdc++.h>
using namespace std;

struct ListNode {
    int val; ListNode* next;
    ListNode(int v) : val(v), next(nullptr) {}
};

ListNode* build(const vector<int>& a) {
    ListNode dummy(0), *t = &dummy;
    for (int x : a) { t->next = new ListNode(x); t = t->next; }
    return dummy.next;
}

ListNode* addTwoNumbers(ListNode* a, ListNode* b) {
    ListNode dummy(0), *t = &dummy;
    int carry = 0;
    while (a || b || carry) {
        int sum = carry;
        if (a) { sum += a->val; a = a->next; }
        if (b) { sum += b->val; b = b->next; }
        carry = sum / 10;
        t->next = new ListNode(sum % 10);
        t = t->next;
    }
    return dummy.next;
}

void print(ListNode* h) { for (; h; h = h->next) cout << h->val << " "; cout << "\n"; }

int main() {
    print(addTwoNumbers(build({2, 4, 3}), build({5, 6, 4})));   // 7 0 8
    return 0;
}
\end{lstlisting}

\begin{complexitybox}
\textbf{Time Complexity.} $O(\max(a,b))$ -- one pass over the longer list.
\textbf{Space Complexity.} $O(\max(a,b))$ for the result list.
\end{complexitybox}

\begin{takeawaybox}
Least-significant-first storage turns list addition into a single forward
carry walk -- the cleanest arithmetic setup. Contrast with Add One's most-
significant-first order, which needs recursion or reversal to reach the
low digit first.
\end{takeawaybox}

\section{Delete All Occurrences of a Key in a Doubly Linked List}
\label{sec:ll-delete-occurrences-dll}

\problemstatement{Given a doubly linked list and a key, delete every node
whose value equals the key, returning the (possibly new) head.}

\begin{patternbox}
The DLL's \texttt{prev} pointer makes this a single-pass sweep: at each
matching node, splice its neighbours together in $O(1)$ using both links,
handling head and tail as the only edge cases.
\end{patternbox}

\subsection*{Splice out each match in place}

Walk the list. When a node's value equals the key, connect its predecessor's
\texttt{next} to its successor and its successor's \texttt{prev} to its
predecessor. If the deleted node was the head, advance the head. Because a
DLL node knows both neighbours, each deletion is constant time with no
re-walk -- the exact advantage over a singly linked list.

\begin{ideabox}[Both Links Make Deletion Local]
Removing a node requires rewiring the links on \emph{both} sides. The
\texttt{prev} pointer supplies the left neighbour instantly, so an arbitrary
node -- including several in one pass -- is removed in $O(1)$ each.
\end{ideabox}

\subsection*{Algorithm}

\begin{enumerate}
  \item Walk from the head. For each node equal to the key: rewire
        \texttt{prev} and \texttt{next} of its neighbours; if it was the
        head, move the head forward.
  \item Continue to the next node.
  \item Return the head.
\end{enumerate}

\subsection*{Dry run}

\dryrun{DLL $1\leftrightarrow2\leftrightarrow2\leftrightarrow3$, key $2$.}

\begin{itemize}
  \item Delete both $2$s by splicing neighbours.
  \item Result $1\leftrightarrow3$.
\end{itemize}

\subsection*{C++ implementation}

\begin{lstlisting}
#include <bits/stdc++.h>
using namespace std;

struct Node {
    int val; Node* prev; Node* next;
    Node(int v) : val(v), prev(nullptr), next(nullptr) {}
};

Node* build(const vector<int>& a) {
    Node dummy(0), *t = &dummy;
    for (int x : a) { Node* n = new Node(x); t->next = n; n->prev = t; t = n; }
    Node* head = dummy.next; if (head) head->prev = nullptr;
    return head;
}

Node* deleteAll(Node* head, int key) {
    Node* cur = head;
    while (cur) {
        Node* nxt = cur->next;
        if (cur->val == key) {
            if (cur->prev) cur->prev->next = cur->next;
            else head = cur->next;                // deleting the head
            if (cur->next) cur->next->prev = cur->prev;
        }
        cur = nxt;
    }
    return head;
}

void print(Node* h) { for (; h; h = h->next) cout << h->val << " "; cout << "\n"; }

int main() {
    print(deleteAll(build({1, 2, 2, 3, 2}), 2));   // 1 3
    return 0;
}
\end{lstlisting}

\begin{complexitybox}
\textbf{Time Complexity.} $O(n)$ -- one pass. \textbf{Space Complexity.}
$O(1)$.
\end{complexitybox}

\begin{takeawaybox}
Deleting all matches in a DLL is a single-pass splice using both links,
with head/tail as the only special cases. This is the everyday DLL deletion
skill that LRU/LFU caches rely on.
\end{takeawaybox}

\section{Find Pairs with a Given Sum in a Sorted DLL}
\label{sec:ll-pairs-sum-dll}

\problemstatement{Given a doubly linked list sorted in ascending order and a
target, find all pairs of nodes whose values sum to the target, using $O(1)$
extra space.}

\begin{patternbox}
A sorted DLL supports the array \textbf{two-pointer} technique directly: one
pointer at the head (smallest), one at the tail (largest). The
\texttt{prev} pointer lets the right pointer move backward, which a singly
linked list cannot do.
\end{patternbox}

\subsection*{Converge from both ends}

Place \texttt{left} at the head and \texttt{right} at the tail. Compare
their sum to the target: if equal, record the pair and move both inward; if
too small, advance \texttt{left} (via \texttt{next}) to increase the sum; if
too large, retreat \texttt{right} (via \texttt{prev}) to decrease it. Stop
when the pointers meet or cross. Sortedness guarantees this monotonic
convergence finds every pair without revisiting.

\begin{ideabox}[The prev Pointer Enables Backward Two Pointers]
The classic two-sum-on-sorted-array needs a right pointer that moves left.
A DLL's \texttt{prev} provides exactly that, so the linear two-pointer sweep
transfers unchanged from arrays to sorted doubly linked lists.
\end{ideabox}

\subsection*{Algorithm}

\begin{enumerate}
  \item \texttt{left = head}; \texttt{right =} last node.
  \item While \texttt{left != right} and not crossed: compare
        \texttt{left->val + right->val} to the target; record on equality
        and move both, else move the appropriate pointer.
  \item Output all recorded pairs.
\end{enumerate}

\subsection*{Dry run}

\dryrun{Sorted DLL $1,2,3,4,6$; target $6$.}

\begin{itemize}
  \item $(1,6)$? $7>6$, move right to $4$. $(1,4)$? $5<6$, move left to
        $2$. $(2,4)=6$ -- record.
  \item Move both; pointers cross. Pair found: $(2,4)$.
\end{itemize}

\subsection*{C++ implementation}

\begin{lstlisting}
#include <bits/stdc++.h>
using namespace std;

struct Node {
    int val; Node* prev; Node* next;
    Node(int v) : val(v), prev(nullptr), next(nullptr) {}
};

Node* build(const vector<int>& a) {
    Node dummy(0), *t = &dummy;
    for (int x : a) { Node* n = new Node(x); t->next = n; n->prev = t; t = n; }
    Node* head = dummy.next; if (head) head->prev = nullptr;
    return head;
}

vector<pair<int,int>> pairsWithSum(Node* head, int target) {
    vector<pair<int,int>> res;
    if (!head) return res;
    Node* left = head;
    Node* right = head;
    while (right->next) right = right->next;       // go to tail
    while (left != right && right->next != left) {
        int sum = left->val + right->val;
        if (sum == target) {
            res.push_back({left->val, right->val});
            left = left->next; right = right->prev;
        } else if (sum < target) left = left->next;
        else right = right->prev;
    }
    return res;
}

int main() {
    for (auto& [a, b] : pairsWithSum(build({1, 2, 3, 4, 6}), 6))
        cout << "(" << a << "," << b << ") ";
    cout << "\n";                                  // (2,4)
    return 0;
}
\end{lstlisting}

\begin{complexitybox}
\textbf{Time Complexity.} $O(n)$ after reaching the tail (also $O(n)$).
\textbf{Space Complexity.} $O(1)$ beyond the output.
\end{complexitybox}

\begin{takeawaybox}
Sorted-DLL pair-sum is the two-pointer two-sum with the right pointer moving
backward via \texttt{prev}. The DLL's backward link is precisely what lets
this array technique work on a list.
\end{takeawaybox}

\section{Remove Duplicates from a Sorted DLL}
\label{sec:ll-remove-duplicates-dll}

\problemstatement{Given a doubly linked list sorted in ascending order,
remove nodes with duplicate values so each value appears once, returning the
head.}

\begin{patternbox}
In a sorted list, duplicates are \textbf{adjacent}, so a single pass
comparing each node to the next -- and splicing out equal successors --
removes all duplicates. The \texttt{prev} pointer keeps the deletions
$O(1)$.
\end{patternbox}

\subsection*{Skip equal neighbours}

Walk from the head. While the current node's value equals its successor's,
unlink the successor (rewire the current node's \texttt{next} and the
node-after's \texttt{prev}). Only advance to the next distinct value once no
more equal successors remain. Sortedness guarantees all copies of a value
are consecutive, so this local check suffices.

\begin{ideabox}[Sorted Means Duplicates Are Adjacent]
Because equal values sit next to each other in a sorted list, comparing
only neighbours catches every duplicate -- no hashing or full scan needed.
The DLL links make each removal constant time.
\end{ideabox}

\subsection*{Algorithm}

\begin{enumerate}
  \item \texttt{cur = head}.
  \item While \texttt{cur} and \texttt{cur->next}: if equal, splice out
        \texttt{cur->next}; else advance \texttt{cur}.
  \item Return the head.
\end{enumerate}

\subsection*{Dry run}

\dryrun{Sorted DLL $1,1,2,3,3$.}

\begin{itemize}
  \item At first $1$: remove the second $1$. At $2$: no dup. At first $3$:
        remove the second $3$.
  \item Result $1,2,3$.
\end{itemize}

\subsection*{C++ implementation}

\begin{lstlisting}
#include <bits/stdc++.h>
using namespace std;

struct Node {
    int val; Node* prev; Node* next;
    Node(int v) : val(v), prev(nullptr), next(nullptr) {}
};

Node* build(const vector<int>& a) {
    Node dummy(0), *t = &dummy;
    for (int x : a) { Node* n = new Node(x); t->next = n; n->prev = t; t = n; }
    Node* head = dummy.next; if (head) head->prev = nullptr;
    return head;
}

Node* removeDuplicates(Node* head) {
    Node* cur = head;
    while (cur && cur->next) {
        if (cur->val == cur->next->val) {
            Node* dup = cur->next;
            cur->next = dup->next;                 // splice out duplicate
            if (dup->next) dup->next->prev = cur;
        } else {
            cur = cur->next;
        }
    }
    return head;
}

void print(Node* h) { for (; h; h = h->next) cout << h->val << " "; cout << "\n"; }

int main() {
    print(removeDuplicates(build({1, 1, 2, 3, 3})));   // 1 2 3
    return 0;
}
\end{lstlisting}

\begin{complexitybox}
\textbf{Time Complexity.} $O(n)$ -- one pass. \textbf{Space Complexity.}
$O(1)$.
\end{complexitybox}

\begin{takeawaybox}
On a sorted list, duplicates are neighbours, so a neighbour-comparison pass
removes them all. Recognising the sorted invariant turns a potential
hashing problem into an $O(1)$-space local sweep.
\end{takeawaybox}

\section{Reverse Nodes in Groups of K}
\label{sec:ll-reverse-k-group}

\problemstatement{Given a linked list and an integer $k$, reverse the nodes
in consecutive groups of $k$. If the final group has fewer than $k$ nodes,
leave it as is. Return the head.}

\begin{patternbox}
This is the iterative reversal applied \textbf{block by block}: for each
group of $k$ nodes, check that $k$ nodes remain, reverse that block in
place, and stitch it between the already-processed part and the rest. A
dummy node anchors the stitching.
\end{patternbox}

\subsection*{Count, reverse a block, reconnect}

Maintain a pointer to the node before the current group. First verify $k$
nodes remain (walk ahead $k$ steps); if not, stop. Reverse the $k$ nodes
with the standard save--flip--advance loop. Then reconnect: the previous
group's tail points to the new (reversed) block head, and the block's tail
points to the following node. Move the ``previous group tail'' pointer to
the block's tail and repeat.

\begin{ideabox}[Reverse in Place, Then Re-stitch the Seams]
The only complexity beyond plain reversal is bookkeeping at the two seams of
each block -- linking the prior tail to the new head and the new tail to the
remainder. A dummy node removes the special case for the very first block.
\end{ideabox}

\subsection*{Algorithm}

\begin{enumerate}
  \item Dummy before head; \texttt{groupPrev = dummy}.
  \item Loop: check $k$ nodes remain; if not, break. Reverse the next $k$
        nodes; reconnect the seams; advance \texttt{groupPrev} to the
        block's (new) tail.
  \item Return \texttt{dummy.next}.
\end{enumerate}

\subsection*{Dry run}

\dryrun{$1\to2\to3\to4\to5$, $k=2$.}

\begin{itemize}
  \item Reverse $[1,2]\to2\to1$; reverse $[3,4]\to4\to3$; $[5]$ left as is.
  \item Result $2\to1\to4\to3\to5$.
\end{itemize}

\subsection*{C++ implementation}

\begin{lstlisting}
#include <bits/stdc++.h>
using namespace std;

struct ListNode {
    int val; ListNode* next;
    ListNode(int v) : val(v), next(nullptr) {}
};

ListNode* build(const vector<int>& a) {
    ListNode dummy(0), *t = &dummy;
    for (int x : a) { t->next = new ListNode(x); t = t->next; }
    return dummy.next;
}

ListNode* reverseKGroup(ListNode* head, int k) {
    ListNode dummy(0); dummy.next = head;
    ListNode* groupPrev = &dummy;

    while (true) {
        ListNode* kth = groupPrev;
        for (int i = 0; i < k && kth; i++) kth = kth->next;   // k-th node
        if (!kth) break;                                       // fewer than k left

        ListNode* groupNext = kth->next;
        ListNode* prev = groupNext;
        ListNode* curr = groupPrev->next;
        while (curr != groupNext) {                            // reverse the block
            ListNode* nx = curr->next;
            curr->next = prev;
            prev = curr;
            curr = nx;
        }
        ListNode* newTail = groupPrev->next;                   // old head -> new tail
        groupPrev->next = kth;                                 // stitch to new head
        groupPrev = newTail;
    }
    return dummy.next;
}

void print(ListNode* h) { for (; h; h = h->next) cout << h->val << " "; cout << "\n"; }

int main() {
    print(reverseKGroup(build({1, 2, 3, 4, 5}), 2));   // 2 1 4 3 5
    return 0;
}
\end{lstlisting}

\begin{complexitybox}
\textbf{Time Complexity.} $O(n)$ -- each node is reversed once.
\textbf{Space Complexity.} $O(1)$.
\end{complexitybox}

\begin{takeawaybox}
Reverse-in-$k$-groups is block-wise iterative reversal with careful seam
stitching. The ``check $k$ remain, reverse, reconnect'' loop is the hard-
problem culmination of the plain reversal skill built earlier.
\end{takeawaybox}

\section{Rotate a Linked List}
\label{sec:ll-rotate}

\problemstatement{Given the head of a linked list and an integer $k$, rotate
the list to the right by $k$ places and return the new head.}

\begin{patternbox}
Rotating right by $k$ is a \textbf{relink}, not a shift: make the list
circular, then break it at the right spot. Since rotating by the length
gives the same list, first reduce $k$ modulo the length.
\end{patternbox}

\subsection*{Circularise, then cut}

Walk to the tail, counting the length $n$, and connect the tail to the head
to form a ring. Reduce $k$ to $k \bmod n$. The new tail is the node at
position $n-k$ from the start; the node after it is the new head. Walk to
that new tail, set its \texttt{next} to null (breaking the ring), and return
the new head. This performs any rotation with a single relink, avoiding
$k$ separate moves.

\begin{ideabox}[Rotation Is One Cut on a Ring]
Joining the list into a circle makes every rotation just a choice of where
to break it. Reducing $k$ modulo $n$ handles $k\ge n$, and the break point
$n-k$ places the last $k$ nodes at the front.
\end{ideabox}

\subsection*{Algorithm}

\begin{enumerate}
  \item Find the length $n$ and the tail; connect tail to head (ring).
  \item $k \gets k \bmod n$; the new tail is $n-k$ steps from the head.
  \item Break after the new tail; return the node following it.
\end{enumerate}

\subsection*{Dry run}

\dryrun{$1\to2\to3\to4\to5$, $k=2$.}

\begin{itemize}
  \item $n=5$, $k=2$; new tail at position $5-2=3$ (node $3$).
  \item Break after $3$; new head $4$. Result $4\to5\to1\to2\to3$.
\end{itemize}

\subsection*{C++ implementation}

\begin{lstlisting}
#include <bits/stdc++.h>
using namespace std;

struct ListNode {
    int val; ListNode* next;
    ListNode(int v) : val(v), next(nullptr) {}
};

ListNode* build(const vector<int>& a) {
    ListNode dummy(0), *t = &dummy;
    for (int x : a) { t->next = new ListNode(x); t = t->next; }
    return dummy.next;
}

ListNode* rotateRight(ListNode* head, int k) {
    if (!head || !head->next || k == 0) return head;
    int n = 1;
    ListNode* tail = head;
    while (tail->next) { tail = tail->next; n++; }
    tail->next = head;                            // make it circular

    k %= n;
    int stepsToNewTail = n - k;
    ListNode* newTail = head;
    for (int i = 1; i < stepsToNewTail; i++) newTail = newTail->next;
    ListNode* newHead = newTail->next;
    newTail->next = nullptr;                       // break the ring
    return newHead;
}

void print(ListNode* h) { for (; h; h = h->next) cout << h->val << " "; cout << "\n"; }

int main() {
    print(rotateRight(build({1, 2, 3, 4, 5}), 2));   // 4 5 1 2 3
    return 0;
}
\end{lstlisting}

\begin{complexitybox}
\textbf{Time Complexity.} $O(n)$ -- one pass to measure, one to reach the
new tail. \textbf{Space Complexity.} $O(1)$.
\end{complexitybox}

\begin{takeawaybox}
Rotation is circularise-then-cut: one relink instead of $k$ moves, with
$k \bmod n$ collapsing large rotations. Turning a list into a ring to
simplify positional surgery is a handy general trick.
\end{takeawaybox}

\section{Flatten a Linked List}
\label{sec:ll-flatten}

\problemstatement{Each node of a linked list has a \texttt{next} pointer (to
the next sub-list head) and a \texttt{bottom} pointer heading a sorted
vertical sub-list. Flatten the whole structure into a single sorted list
linked entirely by \texttt{bottom} pointers.}

\begin{patternbox}
Every vertical sub-list is already sorted, so flattening is \textbf{repeated
merging}: merge two sorted \texttt{bottom}-lists at a time, folding the
horizontal \texttt{next} chain into one long sorted vertical list. It is the
merge routine of merge sort applied across sub-lists.
\end{patternbox}

\subsection*{Merge sub-lists right to left}

Recurse along \texttt{next} to flatten everything to the right of the head
first, then merge the head's vertical list with that flattened remainder.
The merge compares \texttt{bottom} values and links the smaller each time,
producing one sorted \texttt{bottom}-chain. Processing right-to-left means
each merge combines the current sub-list with an already-fully-flattened
tail.

\begin{ideabox}[Fold the Horizontal Lists into One Vertical Merge]
Because each vertical list is sorted, the whole grid flattens by pairwise
merging along the \texttt{next} direction. The result uses only
\texttt{bottom} pointers, discarding the horizontal structure.
\end{ideabox}

\subsection*{Algorithm}

\begin{enumerate}
  \item If \texttt{head} or \texttt{head->next} is null, return
        \texttt{head}.
  \item Recursively flatten \texttt{head->next}.
  \item Merge \texttt{head}'s vertical list with the flattened remainder by
        \texttt{bottom} value; return the merged head.
\end{enumerate}

\subsection*{Dry run}

\dryrun{Two sub-lists: $3\!\downarrow\!7$ and $1\!\downarrow\!4$ (down =
bottom), linked by \texttt{next}.}

\begin{itemize}
  \item Flatten the second: $1\to4$. Merge with $3\to7$.
  \item Result (bottom-linked) $1\to3\to4\to7$.
\end{itemize}

\subsection*{C++ implementation}

\begin{lstlisting}
#include <bits/stdc++.h>
using namespace std;

struct Node {
    int val;
    Node* next;                                   // horizontal
    Node* bottom;                                 // vertical (sorted)
    Node(int v) : val(v), next(nullptr), bottom(nullptr) {}
};

Node* merge(Node* a, Node* b) {                   // merge two bottom-lists
    Node dummy(0), *t = &dummy;
    while (a && b) {
        if (a->val <= b->val) { t->bottom = a; a = a->bottom; }
        else                  { t->bottom = b; b = b->bottom; }
        t = t->bottom;
    }
    t->bottom = a ? a : b;
    return dummy.bottom;
}

Node* flatten(Node* head) {
    if (!head || !head->next) return head;
    head->next = flatten(head->next);             // flatten the rest first
    head = merge(head, head->next);
    return head;
}

int main() {
    Node* a = new Node(3); a->bottom = new Node(7);
    Node* b = new Node(1); b->bottom = new Node(4);
    a->next = b;
    Node* f = flatten(a);
    for (Node* c = f; c; c = c->bottom) cout << c->val << " ";
    cout << "\n";                                  // 1 3 4 7
    return 0;
}
\end{lstlisting}

\begin{complexitybox}
\textbf{Time Complexity.} $O(N)$ total merge work across all
$N$ nodes (with $O(N\cdot m)$ in the naive per-merge bound, $m$ sub-lists).
\textbf{Space Complexity.} $O(m)$ recursion over the \texttt{next} chain.
\end{complexitybox}

\begin{takeawaybox}
Flattening sorted sub-lists is repeated merging along the horizontal
direction, reusing the merge-sort combine step. Recognising ``already
sorted pieces $\Rightarrow$ merge, don't re-sort'' is the key insight; a
priority queue over sub-list heads is an alternative.
\end{takeawaybox}

\section{Clone a Linked List with Random Pointers}
\label{sec:ll-clone-random}

\problemstatement{Each node has a \texttt{next} pointer and a
\texttt{random} pointer to any node (or null). Produce a deep copy of the
list -- new nodes with correctly wired \texttt{next} and \texttt{random}
pointers -- in $O(1)$ extra space.}

\begin{patternbox}
The clever $O(1)$-space method \textbf{interleaves} each copy right after
its original. Then a copy's \texttt{random} is just the original's random's
\texttt{next}. A final pass unweaves the two lists apart.
\end{patternbox}

\subsection*{Interleave, wire randoms, unweave}

\textbf{Pass 1:} for each original node $x$, create a copy $x'$ and insert it
between $x$ and $x.\texttt{next}$, forming
$x\to x'\to x.\texttt{next}\to\ldots$ \textbf{Pass 2:} set each copy's random
via $x'.\texttt{random}=x.\texttt{random}.\texttt{next}$ -- because every
original is immediately followed by its copy, the copy of any target is one
\texttt{next} away. \textbf{Pass 3:} separate the interleaved list back into
the original and the cloned list, restoring \texttt{next} pointers.

\begin{ideabox}[The Copy Sits Next to Its Original]
Interleaving means ``the clone of node $x$'' is always \texttt{x->next}. So
the random pointer of a clone is \texttt{x->random->next} -- computable
without a hash map, which is what removes the $O(n)$ auxiliary space a map-
based clone would use.
\end{ideabox}

\subsection*{Algorithm}

\begin{enumerate}
  \item Insert each node's copy immediately after it.
  \item For each copy, set \texttt{copy->random = orig->random->next} (guard
        null).
  \item Detach the copies into a separate list, restoring the originals'
        \texttt{next}.
\end{enumerate}

\subsection*{Dry run}

\dryrun{$A\to B$ with $A.\texttt{random}=B$, $B.\texttt{random}=B$.}

\begin{itemize}
  \item Interleave: $A\to A'\to B\to B'$. Set $A'.\texttt{random}=
        A.\texttt{random}.\texttt{next}=B'$; $B'.\texttt{random}=B'$.
  \item Unweave: clone is $A'\to B'$ with correct randoms.
\end{itemize}

\subsection*{C++ implementation}

\begin{lstlisting}
#include <bits/stdc++.h>
using namespace std;

struct Node {
    int val;
    Node* next;
    Node* random;
    Node(int v) : val(v), next(nullptr), random(nullptr) {}
};

Node* copyRandomList(Node* head) {
    if (!head) return nullptr;
    // Pass 1: interleave copies
    for (Node* c = head; c; c = c->next->next) {
        Node* copy = new Node(c->val);
        copy->next = c->next;
        c->next = copy;
    }
    // Pass 2: wire random pointers
    for (Node* c = head; c; c = c->next->next)
        if (c->random) c->next->random = c->random->next;
    // Pass 3: unweave
    Node dummy(0), *t = &dummy;
    for (Node* c = head; c; c = c->next) {
        t->next = c->next;                        // pick the copy
        t = t->next;
        c->next = c->next->next;                  // restore original
    }
    return dummy.next;
}

int main() {
    Node* a = new Node(1);
    Node* b = new Node(2);
    a->next = b;
    a->random = b; b->random = b;
    Node* cl = copyRandomList(a);
    for (Node* c = cl; c; c = c->next)
        cout << c->val << "(rand " << (c->random ? c->random->val : -1) << ") ";
    cout << "\n";                                  // 1(rand 2) 2(rand 2)
    return 0;
}
\end{lstlisting}

\begin{complexitybox}
\textbf{Time Complexity.} $O(n)$ -- three linear passes. \textbf{Space
Complexity.} $O(1)$ extra (versus $O(n)$ for a hash-map clone).
\end{complexitybox}

\begin{takeawaybox}
Cloning with random pointers in $O(1)$ space uses interleaving so each
clone sits beside its original, making random wiring a local
\texttt{->next} lookup. It is the standout example of exploiting the list's
own structure instead of an auxiliary map.
\end{takeawaybox}

